\chapter{The simple classification}
\label{ch:gw-simple}

The potential of the quantised law is small, so the matching theorem of \cref{ch:iid} applies
and the transfer of \cref{ch:chain-simple} turns a matching into a quasi-isometry. The bushy
regime runs the same argument over shapes. The obstructions separate the classes.

\section{The potential of the quantised law}

\begin{lemma}[Chain potential bound]
\label{thm:eta-bound}
\lean{ChainClasses.quantised_eta_le}
\leanok
\uses{thm:quantised-law, def:etaG}
For every $\theta_1\in(0,1)$ there is $D_0=D_0(\theta_1)\geq5$ such that for every integer
$D\geq D_0$ the law $p^{(D)}$ of \cref{thm:quantised-law} satisfies
\[
 \eta_{\mathsf P,5/2}\bigl(p^{(D)}\bigr)\leq16\,\theta_1^{D(D-5/2)}\leq10^{-4}.
\]
\end{lemma}

\begin{proof}
\leanok
\uses{thm:quantised-law, thm:dominance}
The classes form a path graph, so local dominance applies with the one-vertex class carrying
most of the mass. The tail terms are summed against the doubly exponential bound
$p^{(D)}_j\leq\theta_1^{D^j-1}$.
\end{proof}

\begin{proof}[Proof of \cref{thm:twovalue}]
\leanok
\proves{thm:twovalue}
\uses{thm:eta-bound, thm:chain-classes, thm:transfer, thm:matching, thm:level}
Quantise both label fields at scale $D$ by \cref{thm:chain-classes}. Their common law has
potential below $10^{-4}$ by \cref{thm:eta-bound}, so \cref{thm:matching} supplies, with
probability at least $1-16\eta$, an automorphism of $\bbB$ matching the two quantised fields at
every vertex. Matched labels lie in equal or adjacent classes, hence are comparable within
$D^2$ by \cref{thm:level}. The transfer lemma \cref{thm:transfer} converts the relabelled field
into a $(D^2+3)$-quasi-isometry, and \cref{thm:isometry} identifies the relabelled tree with
the original. The rate is the potential bound; the almost sure statement follows by letting
$D\to\infty$.
\end{proof}

\begin{proposition}[Universality in the chain regime]
\label{thm:cross-law}
\lean{ChainClasses.crosslaw_ae_tree, ChainClasses.exists_crosslaw_rate_tree}
\leanok
\uses{thm:chain-coupling, thm:transfer, thm:matching}
Let $\theta,\theta'$ satisfy $\theta_1+\theta_2=1$ and $\theta_1'+\theta_2'=1$ with
$\theta_1,\theta_1'\in(0,1)$, and let $\cT(\omega),\cT'(\omega')$ be independent Galton--Watson
trees with these laws. Then almost surely $\cT(\omega)\simeq\cT'(\omega')$, with a rate at the
constant $\gamma_*D^2+3$, where $\gamma$ is as in \cref{thm:chain-coupling}.
\end{proposition}

\begin{proof}
\leanok
\uses{thm:chain-coupling, thm:eta-bound, thm:transfer}
The coupling of \cref{thm:chain-coupling} presents the second law's classes against the first
with a controlled ratio, so a single matching serves both fields. The transfer lemma is then
applied at a real constant.
\end{proof}

\section{The bushy regime}

\begin{lemma}[Shrinking a shape into a support]
\label{thm:shape-shrink}
\lean{ChainClasses.markedQI_shape_shrink}
\leanok
\uses{def:shape, thm:shape-iid}
Let $\theta'$ be supported on $\set{0,1,2}$ and supercritical with $\theta_0'>0$, with shape
law $\mu'$. Every shape $\sigma$ admits a support fix $\sigma^\circ\in\supp\mu'$ with
$\abs{\sigma^\circ}\leq2\abs\sigma$ and a $2$-marked quasi-isometry to $\sigma$, and a shape
may be shrunk into the support at a controlled scale.
\end{lemma}

\begin{lemma}[Shape potential bound]
\label{thm:shape-eta}
\lean{ChainClasses.shape_eta_final}
\leanok
\uses{def:shape-net, thm:shape-mass, thm:dilution}
There are $c=c(\theta)>0$ and $D_0=D_0(\theta)$ such that for every integer $D\geq D_0$,
\[
 \eta_{G_D,5/2}(\mu_D)\leq e^{-cD^{2}}\leq10^{-4}.
\]
Moreover the class of the one-vertex shape carries $\mu_D$-mass at least
$1-e^{-cD^{2}}\geq\tfrac12$.
\end{lemma}

\begin{proof}
\leanok
\uses{thm:shape-mass, thm:dilution, thm:shape-net, thm:dominance}
The one-vertex class dominates because every shape of size at most $D^2$ collapses onto a
single vertex at scale $D$. The remaining terms are counted by entropy dilution and weighted by
the exponential mass bounds, and local dominance closes the estimate.
\end{proof}

\begin{lemma}[Shape coupling]
\label{thm:shape-coupling}
\lean{ChainClasses.exists_isShapeCoupling_etaS}
\leanok
\uses{thm:shape-eta, thm:shape-shrink}
Let $\theta,\theta'$ be supercritical, supported on $\set{0,1,2}$, with $\theta_0,\theta_0'>0$
and shape laws $\mu,\mu'$. There is $D_1=D_1(\theta,\theta')$ and, for $D\geq D_1$, a coupling
of the two shape laws over a common label graph whose potential is again below $10^{-4}$, whose
matched pairs admit a marked quasi-isometry at a controlled scale.
\end{lemma}

\begin{proposition}[Universality in the bushy regime]
\label{thm:hairy}
\lean{ChainClasses.hairy_ae_tree, ChainClasses.hairy_rate_tree_scale}
\leanok
\uses{thm:shape-coupling, thm:glued-transfer, thm:matching}
Let $\theta,\theta'$ be supercritical, supported on $\set{0,1,2}$, with $\theta_0,\theta_0'>0$,
and let $\cT,\cT'$ be independent Galton--Watson trees with these laws conditioned on infinite
diameter. Then almost surely $\cT\simeq\cT'$, with an explicit rate.
\end{proposition}

\begin{proof}
\leanok
\uses{thm:shape-iid, thm:shape-eta, thm:shape-coupling, thm:matching, thm:glued-transfer}
The shapes are i.i.d.\ by \cref{thm:shape-iid} and their net labels have small potential by
\cref{thm:shape-eta}, so \cref{thm:matching} matches the two shape fields along an automorphism
of $\bbB$. Matched shapes admit a marked quasi-isometry, and the glued transfer
\cref{thm:glued-transfer} assembles these into a quasi-isometry of the two assemblies. Across
two laws the shape coupling \cref{thm:shape-coupling} supplies the common label graph.
\end{proof}

\section{The obstructions}

\begin{lemma}[Stability in a tree]
\label{thm:tree-stability}
\lean{BranchingProcess.exists_between_dist_le}
\leanok
\uses{def:qi}
Let $T,T'$ be trees, let $\psi\colon T\to T'$ be a $D$-quasi-isometry, and let $x,y\in T$.
Every vertex of the geodesic $[\psi(x),\psi(y)]$ lies within distance $D$ of the image under
$\psi$ of the geodesic $[x,y]$.
\end{lemma}

\begin{proposition}[Hairs against lines]
\label{thm:hair-separation}
\lean{BranchingProcess.not_quasiIsometric_of_hair_of_line}
\leanok
\uses{thm:tree-stability}
Let $T,T'$ be trees. If $T$ has hairs of unbounded depth and every vertex of $T'$ lies within
distance $R$ of a line of $T'$, then $T\not\simeq T'$.
\end{proposition}

\begin{proposition}[Three rays against the ray]
\label{thm:three-rays}
\lean{BranchingProcess.not_quasiIsometric_rayGraph}
\leanok
\uses{thm:tree-stability}
Let $T$ be a tree containing three rays that pairwise meet only in their common initial vertex.
Then $T$ is not quasi-isometric to the ray $\bbN_0$.
\end{proposition}

\begin{lemma}[The regimes against the obstructions]
\label{thm:regime-obstructions}
\lean{ChainClasses.unbounded_hairs_ae, ChainClasses.nearLine_inTree}
\leanok
\uses{thm:hair-separation, thm:three-rays}
Let $\theta$ be supercritical and supported on $\set{0,1,2}$, and let $\cT$ be a Galton--Watson
tree with law $\theta$ conditioned on infinite diameter. Almost surely $\cT$ carries three rays
meeting only at their common initial vertex; in the chain regime every vertex lies within
bounded distance of a line; and in the bushy regime $\cT$ has hairs of unbounded depth.
\end{lemma}

\begin{proposition}[Bottleneck obstruction]
\label{thm:bottleneck}
\lean{ChainClasses.not_isQIWith_binary_of_thinBalls}
\leanok
\uses{thm:tree-stability}
Let $\cT$ be a tree such that for every $\ell\in\bbN$ there is a vertex $v_\ell$ whose ball
$B_{\cT}(v_\ell,\ell)$ is isometric to a segment of $\bbZ$. Then $\cT$ is not quasi-isometric
to $\bbB$.
\end{proposition}

\begin{theorem}[Separation from the binary tree and the ray]
\label{thm:converse}
\lean{ChainClasses.converse_ae}
\leanok
\uses{thm:bottleneck, thm:three-rays, thm:regime-obstructions}
Let $\theta$ satisfy $\theta_1+\theta_2=1$ with $0<\theta_1<1$, and let $\cT$ be a
Galton--Watson tree with law $\theta$. Then almost surely $\cT$ is quasi-isometric neither to
the binary tree nor to the ray $\bbN_0$.
\end{theorem}

\begin{proof}
\leanok
\uses{thm:bottleneck, thm:three-rays}
Long chains give balls isometric to segments of $\bbZ$ at every scale, which excludes the
binary tree by \cref{thm:bottleneck}. Branch points give three rays meeting only at their
common initial vertex, which excludes the ray by \cref{thm:three-rays}.
\end{proof}
